\documentclass[]{book}
\usepackage{lmodern}
\usepackage{amssymb,amsmath}
\usepackage{ifxetex,ifluatex}
\usepackage{fixltx2e} % provides \textsubscript
\ifnum 0\ifxetex 1\fi\ifluatex 1\fi=0 % if pdftex
  \usepackage[T1]{fontenc}
  \usepackage[utf8]{inputenc}
\else % if luatex or xelatex
  \ifxetex
    \usepackage{mathspec}
  \else
    \usepackage{fontspec}
  \fi
  \defaultfontfeatures{Ligatures=TeX,Scale=MatchLowercase}
\fi
% use upquote if available, for straight quotes in verbatim environments
\IfFileExists{upquote.sty}{\usepackage{upquote}}{}
% use microtype if available
\IfFileExists{microtype.sty}{%
\usepackage[]{microtype}
\UseMicrotypeSet[protrusion]{basicmath} % disable protrusion for tt fonts
}{}
\PassOptionsToPackage{hyphens}{url} % url is loaded by hyperref
\usepackage[unicode=true]{hyperref}
\hypersetup{
            pdftitle={Apuntes para el Taller de Matemáticas - CHyCS},
            pdfauthor={Carlos Ernesto Martinez},
            pdfborder={0 0 0},
            breaklinks=true}
\urlstyle{same}  % don't use monospace font for urls
\usepackage{graphicx,grffile}
\makeatletter
\def\maxwidth{\ifdim\Gin@nat@width>\linewidth\linewidth\else\Gin@nat@width\fi}
\def\maxheight{\ifdim\Gin@nat@height>\textheight\textheight\else\Gin@nat@height\fi}
\makeatother
% Scale images if necessary, so that they will not overflow the page
% margins by default, and it is still possible to overwrite the defaults
% using explicit options in \includegraphics[width, height, ...]{}
\setkeys{Gin}{width=\maxwidth,height=\maxheight,keepaspectratio}
\IfFileExists{parskip.sty}{%
\usepackage{parskip}
}{% else
\setlength{\parindent}{0pt}
\setlength{\parskip}{6pt plus 2pt minus 1pt}
}
\setlength{\emergencystretch}{3em}  % prevent overfull lines
\providecommand{\tightlist}{%
  \setlength{\itemsep}{0pt}\setlength{\parskip}{0pt}}
\setcounter{secnumdepth}{5}
% Redefines (sub)paragraphs to behave more like sections
\ifx\paragraph\undefined\else
\let\oldparagraph\paragraph
\renewcommand{\paragraph}[1]{\oldparagraph{#1}\mbox{}}
\fi
\ifx\subparagraph\undefined\else
\let\oldsubparagraph\subparagraph
\renewcommand{\subparagraph}[1]{\oldsubparagraph{#1}\mbox{}}
\fi

% set default figure placement to htbp
\makeatletter
\def\fps@figure{htbp}
\makeatother

\usepackage[spanish]{babel}
\usepackage{multicol}
\usepackage{titlesec}
\renewcommand{\chaptername}{Capítulo}
\usepackage{fancyhdr}
\usepackage[document]{ragged2e} \pagestyle{fancy} \fancyhf{} \fancyfoot[C]{\thepage} \fancyfoot[L]{\ifodd\value{page}\textit{\small Notas para el Taller de Matemáticas}\else\textit{\small Carlos E Martínez-Rodríguez}\fi} \fancyfoot[R]{\ifodd\value{page}\textit{\small Carlos E Martínez-Rodríguez}\else\textit{\small Notas para el Taller de Matemáticas}\fi}
\usepackage[left=2cm,right=2cm,top=2cm,bottom=2cm]{geometry}
\renewcommand{\chaptermark}[1]{\markboth{#1}{}}
\renewcommand{\sectionmark}[1]{\markright{#1}}
\fancyhead[LE,RO]{\nouppercase{\rightmark}}
\fancyhead[LO,RE]{\nouppercase{\leftmark}}
\newtheorem{Def}{Definición}[chapter]
\newtheorem{Ejem}{Ejemplo}[chapter]
\newtheorem{Ejer}{Ejercicio}[chapter]
\newtheorem{Prop}{Propiedad}[chapter]
\newtheorem{Nota}{Nota}[chapter]
\usepackage{makeidx}
\makeindex

\title{Apuntes para el Taller de Matemáticas - CHyCS}
\author{Carlos Ernesto Martinez}
\date{2023-08-12}

\begin{document}
\maketitle

{
\setcounter{tocdepth}{1}
\tableofcontents
}
\justify

\newpage

\thispagestyle{empty}

\chapter{Presentación}\label{presentaciuxf3n}

En esta página encontrarás ejercicios adicionales a los vistos en clase

\section{Información personal y redes
sociales}\label{informaciuxf3n-personal-y-redes-sociales}

\begin{itemize}
\tightlist
\item
  \href{https://directorioprofesores.uacm.edu.mx/profesor.html?key=2003080144}{Página
  personal} - \href{https://twitter.com/carlosmroder}{Twitter} -
  \href{https://www.facebook.com/TTodoMundo/}{Facebook} -
  \href{https://www.linkedin.com/in/carlos-martinez-r/}{Linkedin}
\end{itemize}

\begin{figure}
\centering
\includegraphics[width=2.08333in]{codigo_qr.png}
\caption{Acceder al documento}
\end{figure}

\chapter{Programa del Taller de
Matemáticas}\label{programa-del-taller-de-matemuxe1ticas}

\section{Propósitos generales y
metodología}\label{propuxf3sitos-generales-y-metodologuxeda}

Este taller se enmarca dentro de los propósitos generales del Programa
de Integración: apoyar y fortalecer la formación del estudiante
proporcionándole el espacio para desarrollar habilidades y herramientas
necesarias para cursar estudios superiores e integrarlo al proyecto
educativo de la UACM.

Se pretende que el estudiante valore a las Matemáticas como un lenguaje
y una herramienta útiles para el análisis de las Ciencias Sociales. De
esta manera, este taller deberá sentar las bases y servir de guía para
que el estudiante pueda continuar con el estudio de las Matemáticas de
acuerdo a las necesidades que surjan en su trayectoria académica.
Además, este taller tiene el propósito de repasar o, en su caso,
introducir contenidos matemáticos básicos y dotar al estudiante de un
espacio en el que pueda continuar con el desarrollo del pensamiento
lógico y abstracto.

El desarrollo de las Matemáticas es una actividad humana y, por lo
tanto, está sujeto a todas las características de una construcción
social. Para este taller se propone una metodología en la que se
desarrollen los contenidos desde una perspectiva que enfatice la
construcción de los conceptos, por ejemplo, desde una aproximación
histórica y/o desde la aplicación de una metodología de trabajo en
equipo. Como resultado de esta perspectiva, se espera que los
estudiantes desarrollen o, en su caso, refuercen las habilidades
necesarias para poder comunicar tanto de manera verbal como escrita sus
razonamientos y argumentaciones. En otras palabras, se busca dar una
perspectiva que no conciba a las Matemáticas como una disciplina
cerrada, sino en continua transformación como realmente lo es.

\section{TEMARIO DEL TALLER}\label{temario-del-taller}

\begin{enumerate}
\def\labelenumi{\arabic{enumi}.}
\tightlist
\item
  Introducción al pensamiento matemático
\end{enumerate}

\textbf{Propósito:} Que el estudiante identifique habilidades de
razonamiento lógico matemático tales como clasificación, reconocimiento
de patrones, generalización, abstracción y deducción para que las valore
como elementos inherentes de las Matemáticas. Que el estudiante
reflexione acerca de las Matemáticas como una actividad humana y conozca
y entienda algunos aspectos de su desarrollo histórico.

\begin{verbatim}
1.1 Razonamiento matemático

  1.1.1 Investigación de patrones numéricos y/o geométricos

  1.1.2 Razonamientos deductivos e inductivos

1.2 Ejemplificar el desarrollo histórico de algún aspecto de las Matemáticas. Sugerencias: Sistemas de numeración, algoritmos para operaciones o demostraciones sencillas con regla y compás.
\end{verbatim}

\begin{enumerate}
\def\labelenumi{\arabic{enumi}.}
\setcounter{enumi}{1}
\tightlist
\item
  Aritmética
\end{enumerate}

\textbf{Propósito:} Que el estudiante reconozca las características
básicas de los diferentes tipos de números reales y su representación en
la recta numérica para que entienda sus propiedades y equivalencias y
las aplique al realizar operaciones básicas entre ellos.

\begin{verbatim}
2.1 Números reales

  2.1.1 Definición de los conjuntos de números naturales, enteros, racionales, reales.

  2.1.2 Ubicación en la recta numérica de números naturales, enteros, racionales y de números expresados en forma decimal. Orden de los números reales.

  2.1.3 Números racionales: significado del numerador y denominador de una fracción en términos de partes en que se divide una unidad y partes de ella que se toman. Fracciones equivalentes.

  2.1.4 Conversión de fracciones en decimales y viceversa

2.2 Operaciones con números reales

  2.2.1 Algoritmos de suma, resta, multiplicación y división de números expresados en forma decimal

  2.2.2 Suma, resta, multiplicación y división de números enteros y jerarquía de las operaciones 2.2.3 Suma, resta, multiplicación y división de números racionales
\end{verbatim}

\begin{enumerate}
\def\labelenumi{\arabic{enumi}.}
\setcounter{enumi}{2}
\tightlist
\item
  Razones y proporciones
\end{enumerate}

\textbf{Propósito:}: Que el estudiante pueda identificar de manera
operativa cuándo dos cantidades son directamente proporcionales y
reconozca el razonamiento que está detrás de la ``regla de tres'' para
que pueda resolver problemas de razones y proporciones enmarcados en el
contexto de las Ciencias Sociales.

\begin{verbatim}
3.1 Razonamiento proporcional

  3.1.1 Interpretaciones de la división en términos de repartir en partes iguales y número de veces que cabe una cantidad en otra. Interpretación de la multiplicación como una suma abreviada. Aplicaciones en problemas de cantidades directamente proporcionales.

  3.1.2 Cantidades unitarias y razón de proporcionalidad directa entre dos variables.

  3.1.3 Cantidades inversamente proporcionales

  3.1.4 Regla de tres simple directa e inversa

  3.1.5 Ejemplos de aplicación en las Ciencias Sociales

3.2 Porcentajes

  3.2.1 El porcentaje como una aplicación del razonamiento proporcional

  3.2.2 Conversiones entre decimales y porcentajes

  3.2.3 Ejemplos de aplicación en las Ciencias Sociales
\end{verbatim}

\begin{enumerate}
\def\labelenumi{\arabic{enumi}.}
\setcounter{enumi}{3}
\tightlist
\item
  Conjuntos
\end{enumerate}

\textbf{Propósito:} Que el estudiante pueda clasificar objetos y
situaciones agrupándolas en conjuntos y asociando las relaciones entre
ellas con las operaciones de conjuntos para que pueda aplicarlas al
análisis de encuestas.

\begin{verbatim}
4.1 Conceptos básicos de la Teoría de Conjuntos

  4.1.1 Conjunto y elemento. Notación de pertenencia. Descripción de conjuntos.

  4.1.2 Igualdad de conjuntos

  4.1.3 Conjunto vacío y conjunto universo. Complemento.

  4.1.4 Subconjunto, subconjunto propio y su notación.

  4.1.5 Diagramas de Venn

4.2 Operaciones entre conjuntos

  4.2.1 Intersección, unión y diferencia. Definiciones.

  4.2.2 Aplicación en conjuntos discretos finitos y usando diagramas de Venn

4.3 Aplicaciones al análisis de encuestas
\end{verbatim}

\begin{enumerate}
\def\labelenumi{\arabic{enumi}.}
\setcounter{enumi}{4}
\tightlist
\item
  Análisis gráfico de la información
\end{enumerate}

\textbf{Propósito:} Que el estudiante conozca las características de la
representación gráfica para que sea capaz de decodificarla y representar
información usando una gráfica.

\begin{verbatim}
5.1 Representación gráfica de dos variables correlacionadas

  5.1.1 Plano cartesiano. Ejes y escalas

  5.1.2 Localización e interpretación de puntos en el plano.

  5.1.3 Interpretación y representación de incrementos verticales y horizontales

5.2 Relación entre variables representada por una gráfica continua en el plano

  5.2.1 Rectas como representación de una relación proporcional. Pendiente y su interpretación como razón de cambio

  5.2.2 Identificar gráficamente: Intervalos de crecimiento y disminución, máximos y mínimos. Relacionar con el signo de la pendiente y su cambio

  5.2.3 Ejemplos de aplicación en las Ciencias Sociales

5.3 Otras representaciones gráficas

  5.3.1 Gráficas de barras

  5.3.2 Gráficas circulares y de área

  5.3.3 Ejemplos de aplicación en las Ciencias Sociales
\end{verbatim}

\begin{enumerate}
\def\labelenumi{\arabic{enumi}.}
\setcounter{enumi}{5}
\tightlist
\item
  Introducción a la Estadística
\end{enumerate}

\textbf{Propósito:} Que el estudiante conozca los elementos básicos de
la estadística y sus fundamentos para que analice información en
contextos simples relacionados con las Ciencias Sociales.

\begin{verbatim}
6.1 Distribuciones de frecuencia y gráficas

  6.1.2 Datos y distribuciones de frecuencia. Agrupamiento de datos.

  6.1.3 Representación gráfica. Histogramas. Opcional: diagramas de tallos y hojas.

  6.1.4 Ejemplos de aplicación en las Ciencias Sociales

6.2 Medidas de tendencia central

  6.2.1 Media, Mediana y Moda. Interpretaciones y diferencias.

  6.2.2 Criterios para escoger alguna medida de tendencia central al realizar un análisis.

  6.2.3 Ejemplos de aplicación en las Ciencias Sociales

6.3 Medidas de dispersión

  6.3.1 Justificación de la introducción de las medidas de dispersión en el análisis estadístico

  6.3.2 Rango y desviación estándar

  6.3.3 Ejemplos de aplicación en las Ciencias Sociales
\end{verbatim}

\section{Metodología}\label{metodologuxeda}

En general, se recomienda que se privilegie la participación activa de
los estudiantes, por ejemplo, permitiendo que ellos propongan ejemplos
que les sean familiares, provocando discusiones colectivas sobre
resultados o dudas y se promueva el trabajo en equipo. El docente deberá
ser un facilitador en la construcción del conocimiento y para esto se
sugiere que evite los tecnicismos, en la medida de lo posible, y que
introduzca palabras nuevas y definiciones sólo hasta que los elementos
que las componen hayan sido discutidos y analizados previamente. Así
mismo, debe procurar no dar respuestas inmediatas a las dudas de los
estudiantes y guiarlos a través de una línea de razonamiento para que,
si se puede, ellos mismos lleguen a las conclusiones correctas (por
ejemplo, a través de preguntas como ``¿cómo sabes qué\ldots{}?'', ``¿qué
evidencia tienes de\ldots{}?'', ``¿qué pasaría si\ldots{}?'').

\section{Sugerencias de materiales para el
taller}\label{sugerencias-de-materiales-para-el-taller}

Para los temas 1, 4 y 6 se puede usar como guía el libro de Miller
{[}1{]}, donde se proponen investigaciones colaborativas que se pueden
proponer como trabajo en equipo. Para los temas 2 y 3, se pueden tomar
actividades del libro de Barkovich y Chernicoff {[}3{]}, que presentan
una serie de actividades diseñadas expresamente para ser trabajadas en
equipo. Para la primera parte del tema 5 se puede usar Chernicoff
{[}5{]}. La segunda parte del tema 1 se puede trabajar a partir de
lecturas basadas en los siguientes libros:

\begin{verbatim}
1. Dioxiadis, Apostolos. El tío Petros y la Conjetura de Goldbach. Madrid: Ediciones B, 2003.

2. Herzenberguer, Hans Magnus. El diablo de los números. Madrid: Siruela, 2003.

3. Infeld, Leopold. El elegido de los dioses. México: Siglo XXI, 2001

4. Mariano Perero, Historia e Historias de las Matemáticas. Grupo Editorial Iberoamérica, 1994.

5. Sagan, Carl. Miles de millones. Madrid: Ediciones B, 2003.

6. Salinger, A. Un matemático lee el periódico. Barcelona: Tusquets Editores, 2002.
\end{verbatim}

\section{Evaluación}\label{evaluaciuxf3n}

Evaluación diagnóstica inicial. Es la evaluación diagnóstica
institucional que presentan los estudiantes de nuevo ingreso. En el
ingreso 2012-II se ha procurado que esta evaluación esté en sintonía con
el presente temario del taller.

Tres evaluaciones formativas parciales.

Se recomienda que para determinar el resultado final del taller se
incluya un examen escrito y un portafolio desarrollado durante el curso
y/o un trabajo de investigación que se puede relacionar con el trabajo
de indagación del taller de ICA.

\section{Bibliografía}\label{bibliografuxeda}

\begin{verbatim}
- MILLER, Charles D., Matemática: Razonamiento y aplicaciones, (1999). 8ª ed. Ed. Addison Wesley. México.

- BARKOVICH, M. y CHERNICOFF, L. Construyendo Matemáticas. Una introducción al razonamiento matemático, (2010). 2ª Ed. Ed. UACM, México. (Disponible sólo en versión electrónica)

- CHERNICOFF, L. Representaciones gráficas, (2010). UACM. Notas para libro. (Disponible sólo en versión electrónica)

- KLINE, Morris. Matemáticas para estudiantes de humanidades, (2002). Ed. Fondo de Cultura Económica. México.
\end{verbatim}

\chapter{El Conjunto de los números
reales}\label{el-conjunto-de-los-nuxfameros-reales}

\section{Semana 1}\label{semana-1}

\begin{enumerate}
\def\labelenumi{\arabic{enumi}.}
\tightlist
\item
  Considerando los siguientes valores de \(n,a\) y \(b\) proporciona las
  coordenadas para cada uno de los puntos \(x_i\), para
  \(i=1,2,\ldots,n\), es decir: \(x_1,x_2,\ldots,x_n\), recuerda que
  \(n\) es el número de subintervalos en los que se divide el segmento
  \(b-a\)
\end{enumerate}

\begin{multicols}{2}
   
    i. n=5, a=2, b=8
   
    ii. n=3, a=1, b=7
   
    iii. n=4, a=2, b=9
   
    iv. n=2, a=1, b=5
   
    v. n=6, a=3, b=7
   
    vi. n=3, a=2, b=6
   
    vii. n=4, a=1, b=6
   
    viii. n=5, a=3, b=9
   
    ix. n=2, a=1, b=4
   
    x. n=4, a=2, b=8
   
    xi. n=5, a=3, b=9
   
    xii. n=2, a=1, b=7
   
    xiii. n=4, a=2, b=6
   
    xiv. n=3, a=1, b=5
   
    xv. n=6, a=4, b=8
\end{multicols}

\begin{enumerate}
\def\labelenumi{\arabic{enumi}.}
\setcounter{enumi}{1}
\item
  Considerando los siguientes valores de \(n,a\) y \(b\) proporciona las
  coordenadas para cada uno de los puntos \(x_i\), para
  \(i=1,2,\ldots,n\), es decir: \(x_1,x_2,\ldots,x_n\), recuerda que
  \(n\) es el número de subintervalos en los que se divide el segmento
  \(b-a\)

  \begin{enumerate}
  \def\labelenumii{\roman{enumii}.}
  \item
    n=5, a=-2, b=8
  \item
    n=3, a=-1, b=7
  \item
    n=4, a=-2, b=9
  \item
    n=2, a=-1, b=5
  \item
    n=6, a=-3, b=7
  \item
    n=3, a=-2, b=6
  \item
    n=4, a=-1, b=6
  \item
    n=5, a=-3, b=9
  \item
    n=2, a=-1, b=4
  \item
    n=4, a=-2, b=8
  \item
    n=5, a=-3, b=9
  \item
    n=2, a=-1, b=7
  \item
    n=4, a=-2, b=6
  \item
    n=3, a=-1, b=5
  \item
    n=6, a=-4, b=8 \textbackslash{}end\{multicols\}
  \end{enumerate}
\item
  Para la siguiente lista de ejercicios determina las coordenadas de los
  puntos \(x_i\), para \(i=1,2,\ldots,n\), es decir:
  \(x_1,x_2,\ldots,x_n\), entre \(-2b\) y \(2b\):
\end{enumerate}

\begin{multicols}{2}
   
    i. n=15, a=22, b=31
   
    ii. n=5, a=9, b=25
   
    iii. n=27, a=7, b=35
   
    iv. n=9, a=16, b=40
   
    v. n=3, a=2, b=15
   
    vi. n=18, a=26, b=9
   
    vii. n=7, a=8, b=28
   
    viii. n=12, a=10, b=39
   
    ix. n=20, a=17, b=34
   
    x. n=10, a=44, b=6
   
    xi. n=25, a=21, b=11
   
    xii. n=8, a=27, b=38
   
    xiii. n=29, a=32, b=43
   
    xiv. n=13, a=19, b=24
   
    xv. n=4, a=37, b=42
\end{multicols}

\section{Semana 2}\label{semana-2}

Para la siguiente lista de valores, determina quién es el mayor de los
dos y ubícalos en la recta real.

\begin{enumerate}
\def\labelenumi{\arabic{enumi}.}
\tightlist
\item
  Determina quién es el mayor de los dos valores y ubícalos en la recta
  real.
\end{enumerate}

\begin{multicols}{2}
  
    i. a = 7/19, b = 11/21
  
   ii. a = 3/13, b = 5/17
  
   iii. a = 9/28, b = 4/12
  
   iv. a = 2/10, b = 6/23
  
   v. a = 14/36, b = 8/40
  
   vi. a = 1/22, b = 15/31
  
   vii. a = 10/29, b = 12/37
  
   viii. a = 5/18, b = 7/32
  
   ix. a = 16/44, b = 20/43
  
   x. a = 6/25, b = 13/35
  
   xi. a = 8/33, b = 3/15
  
   xii. a = 21/42, b = 22/45
  
   xiii. a = 19/41, b = 11/27
  
   xiv. a = 17/34, b = 4/9
  
   xv. a = 23/39, b = 9/20
\end{multicols}

\begin{enumerate}
\def\labelenumi{\arabic{enumi}.}
\setcounter{enumi}{1}
\tightlist
\item
  Determina quién es el mayor de los dos valores y ubícalos en la recta
  real.
\end{enumerate}

\begin{multicols}{2}
   
    i. a = 6/15, b = 18/30
   
   ii. a = 9/18, b = 20/40
   
   iii. a = 7/21, b = 15/45
   
   iv. a = 5/10, b = 10/20
   
   v. a = 3/9, b = 12/36
   
   vi. a = 11/22, b = 24/48
   
   vii. a = 8/24, b = 16/40
   
   viii. a = 4/12, b = 9/27
   
   ix. a = 14/35, b = 28/42
   
   x. a = 2/8, b = 6/24
   
   xi. a = 10/25, b = 21/42
   
   xii. a = 13/26, b = 27/54
   
   xiii. a = 12/30, b = 25/35
   
   xiv. a = 1/5, b = 5/25
   
   xv. a = 16/32, b = 36/45
\end{multicols}

\section{Ejercicios adicionales}\label{ejercicios-adicionales}

\begin{enumerate}
\def\labelenumi{\arabic{enumi}.}
\tightlist
\item
  Considerando los siguientes valores de \(n,a\) y \(b\) proporciona las
  coordenadas para cada uno de los puntos \(x_i\), para
  \(i=1,2,\ldots,n\), es decir: \(x_1,x_2,\ldots,x_n\)
\end{enumerate}

\begin{multicols}{2}
   
    i. n=5, a=2, b=8
   
    ii. n=3, a=1, b=7
   
    iii. n=4, a=2, b=9
   
    iv. n=2, a=1, b=5
   
    v. n=6, a=3, b=7
   
    vi. n=3, a=2, b=6
   
    vii. n=4, a=1, b=6
   
    viii. n=5, a=3, b=9
   
    ix. n=2, a=1, b=4
   
    x. n=4, a=2, b=8
\end{multicols}

\begin{enumerate}
\def\labelenumi{\arabic{enumi}.}
\setcounter{enumi}{1}
\tightlist
\item
  Realiza lo mismo que el ejercicio anterior para las siguientes ternas:
\end{enumerate}

\begin{multicols}{2}
   
    i. n=5, a=-2, b=8
   
    ii. n=3, a=-1, b=7
   
    iii. n=4, a=-2, b=9
   
    iv. n=2, a=-1, b=5
   
    v. n=6, a=-3, b=7
   
    vi. n=3, a=-2, b=6
   
    vii. n=4, a=-1, b=6
   
    viii. n=5, a=-3, b=9
   
    ix. n=2, a=-1, b=4
   
    x. n=4, a=-2, b=8
   
   
    xi. n=5, a=-3, b=9
   
    xii. n=2, a=-1, b=7
   
    xiii. n=4, a=-2, b=6
   
    xiv. n=3, a=-1, b=5
   
    xv. n=6, a=-4, b=8
\end{multicols}

\begin{enumerate}
\def\labelenumi{\arabic{enumi}.}
\setcounter{enumi}{2}
\tightlist
\item
  Para la siguiente lista de ejercicios determina las coordenadas de los
  puntos \(x_i\), para \(i=1,2,\ldots,n\), es decir:
  \(x_1,x_2,\ldots,x_n\), entre \(-2b\) y \(2b\):
\end{enumerate}

\begin{multicols}{2}
   
    i. n=19, a=41, b=33
   
    ii. n=6, a=36, b=45
   
    iii. n=23, a=30, b=14
   
    iv. n=28, a=4, b=18
   
    v. n=17, a=5, b=26
   
    vi. n=11, a=23, b=20
   
    vii. n=2, a=25, b=29
   
    viii. n=21, a=1, b=41
   
    ix. n=14, a=43, b=3
   
    x. n=24, a=12, b=22
\end{multicols}

\begin{enumerate}
\def\labelenumi{\arabic{enumi}.}
\setcounter{enumi}{3}
\tightlist
\item
  Determina cual es el mayor y ubícalos en la recta real
\end{enumerate}

\begin{multicols}{2}
   
    i. a = 24/26, b = 25/30
   
    ii. a = 32/38, b = 7/14
   
    iii. a = 40/44, b = 42/43
   
    iv. a = 1/8, b = 3/11
   
    v. a = 13/37, b = 18/40
   
    vi. a = 1/2, b = 2/3
   
    vii. a = 17/34, b = 38/48
   
    vii. a = 19/38, b = 42/45
   
    ix. a = 22/44, b = 40/43
   
    x. a = 23/46, b = 41/44
\end{multicols}

\chapter{Aritmética}\label{aritmuxe9tica}

\section{Máximo Común Divisor y Mínimo Común
Múltiplo}\label{muxe1ximo-comuxfan-divisor-y-muxednimo-comuxfan-muxfaltiplo}

\begin{enumerate}
\def\labelenumi{\arabic{enumi}.}
\tightlist
\item
  Encuentra el M.C.D de los siguientes pares de números
\end{enumerate}

\begin{multicols}{2}
  
    i. m.c.d(15, 25)
  
    ii. m.c.d(10, 30)
  
    iii. m.c.d(7, 49)
  
    iv. m.c.d(20, 40)
  
    v. m.c.d(12, 48)
  
    vi. m.c.d(9, 18)
  
    vii. m.c.d(4, 44)
  
    viii. m.c.d(21, 35)
  
    ix. m.c.d(28, 42)
  
    x. m.c.d(6, 36)
  
    xi. m.c.d(16, 32)
  
    xii. m.c.d(24, 48)
  
    xiii. m.c.d(8, 16)
  
    xiv. m.c.d(5, 45)
  
    xv. m.c.d(3, 33)
\end{multicols}

\begin{enumerate}
\def\labelenumi{\arabic{enumi}.}
\setcounter{enumi}{1}
\tightlist
\item
  Encuentra el M.C.D de las siguientes ternas de números
\end{enumerate}

\begin{multicols}{2}
  
    i. m.c.d(10, 25, 35)
  
    ii. m.c.d(15, 30, 45)
  
    iii. m.c.d(12, 18, 36)
  
    iv. m.c.d(20, 40, 50)
  
    v. m.c.d(8, 24, 32)
  
    vi. m.c.d(7, 14, 21)
  
    vii. m.c.d(6, 42, 48)
  
    viii. m.c.d(9, 27, 45)
  
    ix. m.c.d(4, 44, 48)
  
    x. m.c.d(16, 32, 48)
  
    xi. m.c.d(22, 33, 44)
  
    xii. m.c.d(25, 35, 49)
  
    xiii. m.c.d(3, 9, 18)
  
    xiv. m.c.d(28, 36, 42)
  
    xv. m.c.d(5, 10, 15)
\end{multicols}

\begin{enumerate}
\def\labelenumi{\arabic{enumi}.}
\setcounter{enumi}{2}
\tightlist
\item
  Encuentra el mínimo común múltipo de los siguientes pares de números
\end{enumerate}

\begin{multicols}{2}
  
    i. m.c.m(10, 25)
  
    ii. m.c.m(15, 30)
  
    iii. m.c.m(12, 18)
  
    iv. m.c.m(20, 40)
  
    v. m.c.m(8, 24)
  
    vi. m.c.m(7, 14)
  
    vii. m.c.m(6, 42)
  
    viii. m.c.m(9, 27)
  
    ix. m.c.m(4, 44)
  
    x. m.c.m(16, 32)
  
    xi. m.c.m(22, 33)
  
    xii. m.c.m(25, 35)
  
    xiii. m.c.m(3, 9)
  
    xiv. m.c.m(28, 36)
  
    xv. m.c.m(5, 10)
\end{multicols}

\begin{enumerate}
\def\labelenumi{\arabic{enumi}.}
\setcounter{enumi}{3}
\tightlist
\item
  Encuentra el mínimo común múltipo de las siguientes ternas de números
\end{enumerate}

\begin{multicols}{2}
  
    i. m.c.m(10, 25, 35)
  
    ii. m.c.m(15, 30, 45)
  
    iii. m.c.m(12, 18, 36)
  
    iv. m.c.m(20, 40, 50)
  
    v. m.c.m(8, 24, 32)
  
    vi. m.c.m(7, 14, 21)
  
    vii. m.c.m(6, 42, 48)
  
    viii. m.c.m(9, 27, 45)
  
    ix. m.c.m(4, 44, 48)
  
    x. m.c.m(16, 32, 48)
  
    xi. m.c.m(22, 33, 44)
  
    xii. m.c.m(25, 35, 49)
  
    xiii. m.c.m(3, 9, 18)
  
    xiv. m.c.m(28, 36, 42)
  
    xv. m.c.m(5, 10, 15)
\end{multicols}

\begin{enumerate}
\def\labelenumi{\arabic{enumi}.}
\setcounter{enumi}{4}
\tightlist
\item
  Encuentra el M.C.D o m.c.m según se indique
\end{enumerate}

\begin{multicols}{2}
  
    i. m.c.m(7, 14)
  
    ii. m.c.m(10, 20)
  
    iii. m.c.m(18, 36)
  
    iv. m.c.m(25, 35)
  
    v. m.c.m(9, 27)
  
    vi. m.c.m(11, 22)
  
    vii. m.c.m(30, 45)
  
    viii. m.c.m(16, 32)
  
    ix. m.c.m(12, 48)
  
    x. m.c.m(6, 18)
  
    xi. m.c.m(4, 8, 16)
  
    xii. m.c.m(6, 12, 18)
  
    xiii. m.c.m(15, 25, 35)
  
    xiv. m.c.m(10, 20, 30)
  
    xv. m.c.m(21, 28, 42)
  
    xvi. m.c.d(6, 18)
  
    xvii. m.c.d(9, 27)
  
    xviii. m.c.d(15, 25)
  
    xix. m.c.d(24, 48)
  
    xx. m.c.d(35, 49)
  
    xxi. m.c.d(4, 8, 12)
  
    xxii. m.c.d(9, 18, 27)
  
    xxiii. m.c.d(7, 14, 21)
  
    xxiv. m.c.d(20, 40, 50)
  
    xxv. m.c.d(16, 32, 48)
\end{multicols}

\section{Operaciones con fracciones: suma, resta, multiplicación y
división}\label{operaciones-con-fracciones-suma-resta-multiplicaciuxf3n-y-divisiuxf3n}

\subsection{Suma de fracciones}\label{suma-de-fracciones}

\begin{enumerate}
\def\labelenumi{\arabic{enumi}.}
\tightlist
\item
  Realiza las siguientes sumas de fracciones
\end{enumerate}

\begin{multicols}{2}
  
    i. $\frac{5}{25} + \frac{8}{40}$
  
    ii. $\frac{6}{18} + \frac{7}{35}$
  
    iii. $\frac{4}{28} + \frac{9}{45}$
  
    iv. $\frac{3}{15} + \frac{6}{36}$
  
    v. $\frac{18}{36} + \frac{10}{52}$
  
    vi. $\frac{7}{21} + \frac{5}{25}$
  
    vii. $\frac{32}{14} + \frac{9}{45}$
  
    viii. $\frac{68}{48} + \frac{64}{24}$
  
    ix. $\frac{36}{45} + \frac{28}{10}$
  
    xi. $\frac{32}{27} + \frac{56}{36}$
  
    x. $\frac{9}{36} + \frac{24}{64}$
  
    xii. $\frac{1}{11} + \frac{8}{40}$
  
    xiii. $\frac{15}{25} + \frac{44}{20}$
  
    xiv. $\frac{10}{50} + \frac{3}{15}$
  
    xv. $\frac{7}{28} + \frac{6}{36}$
  
    xvi. $\frac{42}{16} + \frac{9}{45}$
  
    xvii. $\frac{6}{30} + \frac{5}{25}$
  
    xviii. $\frac{44}{20} + \frac{27}{63}$
  
    xix. $\frac{8}{40} + \frac{1}{11}$
  
    xx. $\frac{42}{18} + \frac{6}{18}$
  
    xxi. $\frac{1}{4} + \frac{2}{3}$
  
    xxii. $\frac{3}{5} + \frac{2}{7}$
  
    xxiii. $\frac{2}{9} + \frac{5}{6}$
  
    xxiv. $\frac{7}{12} + \frac{3}{8}$
  
    xxv. $\frac{4}{9} + \frac{3}{5} + \frac{2}{7}$
  
    xxvi. $\frac{5}{8} + \frac{3}{10} + \frac{2}{5}$
  
    xxvii. $\frac{6}{11} + \frac{4}{9} + \frac{3}{7}$
  
    xxviii. $\frac{7}{10} + \frac{5}{8} + \frac{4}{9}$
  
    xxix. $\frac{2}{3} + \frac{4}{5} + \frac{6}{7} + \frac{8}{9}$
  
    xxx. $\frac{9}{11} + \frac{5}{7} + \frac{3}{4} + \frac{2}{3}$
  
    xxxi. $\frac{2}{5} + \frac{7}{9} + \frac{4}{6} + \frac{1}{8}$
  
    xxxii. $\frac{5}{8} + \frac{2}{3} + \frac{4}{7} + \frac{6}{9}$
  
    xxxiii. $\frac{3}{4} + \frac{4}{5} + \frac{5}{6} + \frac{6}{7} + \frac{7}{8}$
  
    xxxiv. $\frac{8}{9} + \frac{7}{10} + \frac{6}{11} + \frac{5}{12} + \frac{4}{13}$
  
    xxxv. $\frac{1}{2} + \frac{3}{4} + \frac{5}{6} + \frac{7}{8} + \frac{9}{10}$
  
    xxxvi. $\frac{11}{12} + \frac{10}{11} + \frac{9}{10} + \frac{8}{9} + \frac{7}{8}$
\end{multicols}

\subsection{Resta de fracciones}\label{resta-de-fracciones}

\begin{enumerate}
\def\labelenumi{\arabic{enumi}.}
\setcounter{enumi}{1}
\tightlist
\item
  Realiza las siguientes restas de fracciones
\end{enumerate}

\begin{multicols}{2}
 
    i. $\frac{3}{7} - \frac{2}{5}$
 
    ii. $\frac{5}{12} - \frac{1}{6}$
 
    iii. $\frac{2}{9} - \frac{4}{11}$
 
    iv. $\frac{7}{15} - \frac{3}{10}$
 
    v. $\frac{9}{20} - \frac{7}{25}$
 
    vi. $\frac{11}{24} - \frac{5}{18}$
 
    vii. $\frac{3}{14} - \frac{2}{21}$
 
    vii. $\frac{6}{13} - \frac{2}{39}$
 
    ix. $\frac{5}{12} - \frac{3}{8}$
 
    x. $\frac{7}{10} - \frac{4}{25}$
 
    xi. $\frac{19}{22} - \frac{11}{44}$
 
    xii. $\frac{8}{15} - \frac{5}{12}$
 
    xiii. $\frac{14}{27} - \frac{5}{18}$    
 
    xiv. $\frac{21}{35} - \frac{12}{25}$
 
    xv. $\frac{5}{12} - \frac{3}{8}$
 
    xvi. $\frac{11}{20} - \frac{7}{15}$
 
    xvii. $\frac{7}{16} - \frac{3}{10}$
 
    xviii. $\frac{13}{24} - \frac{5}{16}$
 
    xix. $\frac{15}{28} - \frac{9}{35}$
 
    xx. $\frac{17}{36} - \frac{11}{27}$
 
    xxi. $\frac{7}{18} - \frac{5}{12}$
 
    xxii. $\frac{16}{33} - \frac{8}{27}$
 
    xxiii. $\frac{10}{21} - \frac{3}{14}$
 
    xxiv. $\frac{25}{51} - \frac{7}{34}$
 
    xxv. $\frac{17}{38} - \frac{6}{19}$
 
    xxvi. $\frac{3}{4} - \frac{1}{2}$
 
    xxvii. $\frac{5}{7} - \frac{2}{5}$
 
    xxviii. $\frac{7}{9} - \frac{3}{4}$
 
    xxix. $\frac{2}{5} - \frac{1}{3}$
 
    xxx. $\frac{4}{7} - \frac{2}{9} - \frac{1}{3}$
 
    xxxi. $\frac{9}{11} - \frac{3}{8} - \frac{4}{7}$
 
    xxxii. $\frac{5}{6} - \frac{2}{5} - \frac{3}{7}$
 
    xxxiii. $\frac{8}{9} - \frac{4}{7} - \frac{1}{6}$
 
    xxxiv. $\frac{7}{10} - \frac{2}{3} - \frac{3}{8} - \frac{1}{5}$
 
    xxxv. $\frac{11}{12} - \frac{5}{9} - \frac{2}{7} - \frac{3}{11}$
 
    xxxvi. $\frac{4}{5} - \frac{1}{3} - \frac{2}{9} - \frac{5}{8}$
 
    xxxvii. $\frac{9}{10} - \frac{3}{7} - \frac{4}{11} - \frac{1}{8}$
 
    xxxviii. $\frac{5}{6} - \frac{2}{5} - \frac{3}{7} - \frac{4}{9} - \frac{1}{10}$
 
    xxxix. $\frac{7}{8} - \frac{4}{11} - \frac{5}{12} - \frac{3}{7} - \frac{2}{9}$
 
    xl. $\frac{3}{4} - \frac{1}{6} - \frac{5}{8} - \frac{2}{9} - \frac{4}{11}$
 
    xli. $\frac{10}{11} - \frac{3}{8} - \frac{4}{9} - \frac{5}{12} - \frac{1}{13}$
\end{multicols}

\subsection{Multiplicación de
fracciones}\label{multiplicaciuxf3n-de-fracciones}

\begin{enumerate}
\def\labelenumi{\arabic{enumi}.}
\setcounter{enumi}{2}
\tightlist
\item
  Multiplicación de fracciones
\end{enumerate}

\begin{multicols}{2}
 
    i. $\frac{1}{2} \times \frac{3}{4}$
 
    ii. $\frac{2}{5} \times \frac{4}{7}$
 
    iii. $\frac{3}{8} \times \frac{5}{6}$
 
    iv. $\frac{4}{9} \times \frac{2}{3}$
 
    v. $\frac{2}{3} \times \frac{4}{5}$
 
    vi. $\frac{3}{7} \times \frac{5}{8}$
 
    vii. $\frac{4}{9} \times \frac{6}{7}$
 
    viii. $\frac{1}{5} \times \frac{3}{4}$
 
    ix. $\frac{2}{5} \times \frac{3}{7} \times \frac{4}{9}$
 
    x. $\frac{6}{11} \times \frac{4}{9} \times \frac{3}{7}$
 
    xi. $\frac{1}{6} \times \frac{5}{8} \times \frac{2}{4}$
 
    xii.$\frac{3}{8} \times \frac{4}{7} \times \frac{5}{6}$
 
    xiii. $\frac{2}{3} \times \frac{3}{5} \times \frac{4}{7}$
 
    xiv$\frac{5}{8} \times \frac{2}{9} \times \frac{3}{10}$
 
    xv. $\frac{1}{4} \times \frac{7}{12} \times \frac{3}{5}$
 
    xvi. $\frac{3}{5} \times \frac{4}{7} \times \frac{5}{8}$
 
    xvii. $\frac{2}{5} \times \frac{3}{7} \times \frac{4}{9} \times \frac{5}{11}$
 
    xviii. $\frac{7}{9} \times \frac{5}{8} \times \frac{3}{7} \times \frac{4}{11}$
 
    xix. $\frac{1}{3} \times \frac{4}{5} \times \frac{6}{7} \times \frac{8}{9}$
 
    xx. $\frac{2}{7} \times \frac{3}{8} \times \frac{5}{11} \times \frac{7}{12}$
 
    xxi. $\frac{2}{3} \times \frac{3}{5} \times \frac{4}{7} \times \frac{5}{9}$
 
    xxii. $\frac{7}{10} \times \frac{5}{8} \times \frac{3}{7} \times \frac{4}{11}$
 
    xxiii. $\frac{1}{4} \times \frac{7}{9} \times \frac{3}{6} \times \frac{5}{8}$
 
    xxiv. $\frac{2}{7} \times \frac{3}{8} \times \frac{5}{9} \times \frac{7}{10}$
 
    xxv. $\frac{2}{3} \times \frac{3}{5} \times \frac{4}{7} \times \frac{5}{9} \times \frac{6}{11}$
 
    xxvi. $\frac{5}{8} \times \frac{2}{9} \times \frac{3}{10} \times \frac{4}{11} \times \frac{7}{13}$
 
    xxvii. $\frac{1}{4} \times \frac{7}{12} \times \frac{3}{5} \times \frac{5}{8} \times \frac{2}{9}$
 
    xxviii. $\frac{3}{5} \times \frac{4}{7} \times \frac{5}{8} \times \frac{6}{9} \times \frac{7}{10}$
 
    xxix. $\frac{2}{3} \times \frac{3}{5} \times \frac{4}{7} \times \frac{5}{9} \times \frac{6}{11}$
 
    xxx. $\frac{5}{8} \times \frac{2}{9} \times \frac{3}{10} \times \frac{4}{11} \times \frac{7}{12}$
 
    xxxi. $\frac{1}{4} \times \frac{7}{12} \times \frac{3}{5} \times \frac{5}{8} \times \frac{2}{9}$
 
    xxxii. $\frac{3}{5} \times \frac{4}{7} \times \frac{5}{8} \times \frac{6}{9} \times \frac{7}{10}$
\end{multicols}

\subsection{División de fracciones}\label{divisiuxf3n-de-fracciones}

\begin{enumerate}
\def\labelenumi{\arabic{enumi}.}
\setcounter{enumi}{3}
\tightlist
\item
  División de fracciones
\end{enumerate}

\begin{multicols}{2}
 
    i. $\frac{2}{3} \div \frac{4}{5}$
 
    ii. $\frac{3}{7} \div \frac{5}{8}$
 
    iii. $\frac{4}{9} \div \frac{6}{7}$
 
    iv. $\frac{1}{5} \div \frac{3}{4}$
 
    v. $\frac{3}{4} \div \frac{2}{3}$
 
    vi. $\frac{5}{8} \div \frac{3}{7}$
 
    vii. $\frac{6}{7} \div \frac{4}{9}$
 
    viii. $\frac{2}{3} \div \frac{1}{5}$
 
    ix. $\frac{2}{5} \div (\frac{3}{7} \div \frac{4}{9})$
 
    x. $\frac{6}{11} \div (\frac{4}{9} \div \frac{3}{7})$
 
    xi. $\frac{1}{6} \div (\frac{5}{8} \div \frac{2}{4})$
 
    xii. $\frac{3}{8} \div (\frac{4}{7} \div \frac{5}{6})$
 
    xiii. $\frac{4}{9} \div (\frac{3}{5} \div \frac{2}{7})$
 
    xiv. $\frac{5}{8} \div (\frac{3}{10} \div \frac{2}{5})$
 
    xv. $\frac{6}{11} \div (\frac{4}{9} \div \frac{3}{7})$
 
    xvi. $\frac{7}{10} \div (\frac{5}{8} \div \frac{4}{9})$
 
    xvii. $(\frac{2}{3} \div \frac{3}{5}) \div (\frac{4}{7} \div \frac{5}{9})$
 
    xviii. $(\frac{7}{10} \div \frac{5}{8}) \div (\frac{3}{7} \div \frac{4}{11})$
 
    xix. $(\frac{1}{4} \div \frac{7}{9}) \div (\frac{3}{6} \div \frac{5}{8})$
 
    xx. $(\frac{2}{7} \div \frac{3}{8}) \div (\frac{5}{9} \div \frac{7}{10})$
 
    xxi. $(\frac{2}{3} \div \frac{4}{5}) \div (\frac{6}{7} \div \frac{8}{9})$
 
    xxii. $(\frac{9}{11} \div \frac{5}{7}) \div (\frac{3}{4} \div \frac{2}{3})$
 
    xxiii. $(\frac{2}{5} \div \frac{7}{9}) \div (\frac{4}{6} \div \frac{1}{8})$
 
    xxiv. $(\frac{5}{8} \div \frac{2}{3}) \div (\frac{4}{7} \div \frac{6}{9})$
 
    xxv. $\frac{2}{3} \div (\frac{3}{5} \div \frac{4}{7}) \div (\frac{5}{9} \div \frac{6}{11})$
 
    xxvi. $\frac{5}{8} \div (\frac{2}{9} \div \frac{3}{10}) \div (\frac{4}{11} \div \frac{7}{12})$
 
    xxvii. $\frac{1}{4} \div (\frac{7}{12} \div \frac{3}{5}) \div (\frac{5}{8} \div \frac{2}{9})$
 
    xxviii. $\frac{3}{5} \div (\frac{4}{7} \div \frac{5}{8}) \div (\frac{6}{9} \div \frac{7}{10})$
 
    xxix. $\frac{3}{4} \div (\frac{4}{5} \div \frac{5}{6}) \div (\frac{6}{7} \div \frac{7}{8})$
 
    xxx. $\frac{8}{9} \div (\frac{7}{10} \div \frac{6}{11}) \div (\frac{5}{12} \div \frac{4}{13})$
 
    xxxi. $\frac{1}{2} \div (\frac{3}{4} \div \frac{5}{6}) \div (\frac{7}{8} \div \frac{9}{10})$
 
    xxxii. $\frac{11}{12} \div (\frac{10}{11} \div \frac{9}{10}) (\div \frac{8}{9} \div \frac{7}{8})$
\end{multicols}

\chapter{Temas adicionales}\label{temas-adicionales}

\section{A. Biografías de Científicas Mexicanas y de la
Universidad}\label{a.-biografuxedas-de-cientuxedficas-mexicanas-y-de-la-universidad}

A continuación les presento una lista con algunas de las científicas
mexicanas y de nuestra universidad de las cuales se tendrá que elaborar
una biografía.

A. Nacionales

\begin{multicols}{2} 

      1. Silvia Ruíz Acosta
  
      2. Ana Meda Guardiola
  
      3. María Emilia Caballero
  
      4. Patricia Saavedra Barrera
  
      5. Eliane Rodríguez
  
      6. Gloria Soberón Chávez
  
      7. Blanca Haydé Ruiz Ordaz
  
      8. Leticia Rocha Zavaleta
  
      9. Rosa Laura Camarena Mejía
  
      10. María Elena Flores Carrasco
  
      11. Atocha Aliseda Llera
  
      12. María De Los Ángeles Eraña Lagos
  
      13. Vivette García Deister
  
      14. Siobhan Fenella Guerrero Mc Manus
  
      15. María Teresa Muñoz Sánchez
\end{multicols}

B. Universidad - Ciencia y Tecnología

\begin{multicols}{2}

      1. Elisa Azuara Liceaga
  
      2. Yesica A. Peña Castañeda
  
      3. Martha Yocupiccio
  
      4. Elizbeth Álvarez Sánchez
  
      5. Minerva Camacho Nuez
  
      6. Lilia López Casanovas
  
      7. Selene Zárate Guerra
  
      8. Rosa Margarita Álvarez González 
  
      9. Gloria Elena Londoño Mejía
  
      10. María Claudia Roldán Ahumada
\end{multicols}

C. Universidad - Ciencias Sociales y Humanidades

\begin{multicols}{2}

      1. Martha Rizo
  
      2. Tania Rodríguez
  
      3. María Rodríguez
  
      4. Erika Araiza Díaz
  
      5. Mariana Berlanga Gayón
  
      6. Pilar Calveiro Garrido
  
      7. Silvia Castellanos Castro
  
      8. Hilda Rosina Conde Zambada
  
      9. Carmen Ross
  
      10. Rodríguez Sánchez María del Carmen
\end{multicols}

\section{B. Línea de tiempo de la
Matemática}\label{b.-luxednea-de-tiempo-de-la-matemuxe1tica}

En este trabajo el objetivo es revisar el desarrollo de las matemáticas
a lo largo de la historia, desde los Griegos hasta la edad media, luego
de la edad media hasta la aparición de Galileo Galilei, Isaac Newton,
Johannes Keppler entre otros para posteriormente el desarrollo de la
ciencia hasta la aparición de Albert Einstein. A continuación la
evolución de las matemáticas hasta Stephen Hawking, y cerrando la
revisión con los trabajos de Gregory Perelman y hasta Maryam Mirzakhani.

Ojalá y la revisión de estas grandes personalidades sirva como incentivo
en su deseo de conocer aún más a las Matemáticas y continuar su estudio
hasta lograr un dominio satisfactorio de las mismas.

\section{C. Las matemáticas en el
cine}\label{c.-las-matemuxe1ticas-en-el-cine}

Las matemáticas han tenido un papel central en algunas películas, por lo
tanto es importante que conozcamos algunas de ellas que han sido un
pilar en la difusión y en algunos casos romantización/dramatización de
sus personajes principales.

Las películas de las cuales es preciso que entreguen una reseña son las
siguientes:

\begin{multicols}{2}

      1. Donald en el país de las matemáticas
  
      2. Nuestro amigo el átomo
  
      3. 21 Blackjack
  
      4. Una mente brillante
  
      5. Mente indomable
  
      6. Ramanujan: el hombre que conocía el infinito
  
      7. Código Enigma
  
      8. La teoría del todo
\end{multicols}

\end{document}
